\section{Number Theory}
\label{sec:numbertheory}

The previous section computed \emph{numbers}: quantities carried to whatever precision
you asked for, where the last digit is always an approximation and the question is only
how good an approximation. This section returns to the opposite pole---the exact
arithmetic of the whole numbers, where nothing rounds and nothing overflows. An integer
in Mathilda is not a 64-bit machine word straining against its limits but, the moment it
outgrows one, a GMP\index{GMP} big integer that grows without bound, so
\(\B{Fibonacci}[200]\) and \(2^{1000}\) are computed to the last digit as readily as
\(2+2\). Number theory is the mathematics of exactly these objects, and it is where a
computer algebra system shows a face very unlike a calculator's: not \emph{how much} but
\emph{how}---how an integer factors, which integers are prime, how they behave under a
modulus, how a rational unfolds into a continued fraction.

We take the subject in the classical order. First the greatest common divisor and the
Euclidean algorithm that computes it (\B{GCD}, \B{LCM}, \B{ExtendedGCD}, \B{CoprimeQ},
\B{Divisible}); then primes and how Mathilda tests, counts, and finds them (\B{PrimeQ},
\B{Prime}, \B{PrimePi}, \B{NextPrime}, \mcode{RandomPrime}, \mcode{PrimePowerQ}); then the
factorization pipeline that breaks a composite into primes (\B{FactorInteger},
\B{IntegerExponent}, \B{SquareFreeQ}), together with a special-purpose factoring method,
rational base descent, that was designed by Mathilda's author; the
classical arithmetic functions built on that factorization (\B{Divisors},
\B{DivisorSigma}, \B{EulerPhi}, \B{MoebiusMu}, \B{LiouvilleLambda}, \B{PrimeOmega},
\B{PrimeNu}); modular arithmetic and the structure of the residues (\B{PowerMod},
\B{MultiplicativeOrder}, \B{PrimitiveRoot}, \B{JacobiSymbol}); continued fractions
(\B{ContinuedFraction}, \B{FromContinuedFraction}); the digit and combinatorial functions
(\B{IntegerDigits}, \B{Binomial}, \B{PartitionsP}); and we close by assembling almost all
of it into a working RSA cryptosystem.

\subsection{The Euclidean algorithm and B\'ezout}
\label{ssec:nt-euclid}

The greatest common divisor is the foundation stone. \B{GCD} returns the largest integer
dividing all of its arguments and \B{LCM}\index{least common multiple} the smallest they all divide; both take any
number of arguments and both were met in passing in \S\ref{sec:arithmetic}. What makes
them a starting point rather than a footnote is the algorithm underneath---Euclid's, the
oldest nontrivial algorithm still in daily use\index{Euclidean algorithm}---and the
identity that ties the two functions together:

\begin{codepairs}
\pair{number-theory/euclid/1}{The largest integer dividing both \(24\) and \(36\).}
\pair{number-theory/euclid/2}{And the smallest that \(4\) and \(6\) both divide.}
\pair{number-theory/euclid/3}{Their product \(\gcd\cdot\operatorname{lcm}\) is formed.}
\pair{number-theory/euclid/4}{It equals the product of the numbers themselves---a
classic identity, true for any pair.}
\end{codepairs}

\B{ExtendedGCD} does more than report the gcd: it returns a \emph{B\'ezout
certificate}\index{Bezout's identity@B\'ezout's identity}, integers expressing the gcd as
a linear combination of the inputs. For \(\B{ExtendedGCD}[a,b]=\{g,\{s,t\}\}\) you always
have \(sa+tb=g\):

\begin{codepairs}
\pair{number-theory/euclid/5}{Here \(\gcd(240,46)=2\), with the certificate coefficients
\(s=-9\), \(t=47\).}
\pair{number-theory/euclid/6}{Read in the right order they check out: \((-9)\cdot240 +
47\cdot46 = 2\). These coefficients are exactly what inverts one number modulo another---
the engine behind the modular division of \S\ref{ssec:nt-modular}.}
\pair{number-theory/euclid/7}{\B{CoprimeQ}\index{coprime} asks the yes/no question directly: \(14\) and
\(15\) share no factor.}
\pair{number-theory/euclid/8}{\(14\) and \(21\) share a \(7\), so they are not coprime.}
\pair{number-theory/euclid/9}{\B{Divisible} tests one-way divisibility: \(4\mid100\).}
\pair{number-theory/euclid/10}{\(7\) does not divide \(100\).}
\end{codepairs}

\usagebox{ExtendedGCD}

\calloutnote{Under the hood}{Euclid's algorithm replaces the pair \((a,b)\) with
\((b,\,a\bmod b)\) until the second entry is zero; the extended form carries the two
B\'ezout coefficients along the same recurrence. Mathilda does not run this in
interpreted code---it hands the integers to GMP's \texttt{mpz\_gcd} and
\texttt{mpz\_gcdext}, whose subquadratic ``half-GCD'' handles thousand-digit inputs in
microseconds. The partial quotients Euclid discards on the way are, as
\S\ref{ssec:nt-contfrac} will show, precisely a continued fraction~\cite{knuth-taocp2}.}

Euclid's algorithm needs only a division with a remainder that is \emph{smaller} than the
divisor, and the Gaussian integers\index{Gaussian integers!Euclidean algorithm}
\(\mathbb{Z}[i]=\{a+bi : a,b\in\mathbb{Z}\}\) have one: divide in \(\mathbb{Q}(i)\), round
both parts to the nearest integer, and the remainder's norm \(a^{2}+b^{2}\) is at most half
the divisor's. So the whole of this subsection carries over, and Mathilda's divisibility
functions accept Gaussian arguments directly. Because \(\pm1\) and \(\pm i\) are all units,
a gcd is defined only up to a unit; Mathilda reports the associate in the first quadrant.

\begin{codepairs}
\pair{number-theory/gaussian/1}{\(5=(2+i)(2-i)\) and \(3+4i=(2+i)^{2}\) share the Gaussian
prime \(2+i\)---and only once, since \(2-i\) does not divide \(3+4i\).}
\pair{number-theory/gaussian/2}{\(3+i=(1+i)(2-i)\) divides \(10\) outright.}
\pair{number-theory/gaussian/3}{Conjugate primes are coprime, so their lcm is their product.}
\pair{number-theory/gaussian/4}{A B\'ezout certificate in \(\mathbb{Z}[i]\): \(0\cdot(3+i) +
i\,(2-i) = 1+2i\), the first-quadrant associate of the common factor \(2-i\).}
\pair{number-theory/gaussian/5}{The remainder \(1+i\) has norm \(2\), below the divisor's
\(13\).}
\pair{number-theory/gaussian/6}{Modular powers too: \((1+i)^{10}=32i\equiv-1+i\).}
\end{codepairs}

Two edge conventions are worth knowing. With a single argument, \mcode{CoprimeQ} asks
whether that number is coprime to \emph{itself}---true only of the units---and a modulus of
zero has no residues at all:

\begin{codepairs}
\pair{number-theory/gaussian/7}{\(\gcd(3,3)=3\), so \(3\) is not ``coprime''.}
\pair{number-theory/gaussian/8}{\(-1\) is a unit, so it is.}
\pair{number-theory/gaussian/9}{\(7\bmod0\) is undefined; \B{Mod} says so, with an
\mcode{Mod::indet} message on the side.}
\end{codepairs}

\subsection{Primes and primality}
\label{ssec:nt-primes}

A prime is an integer greater than \(1\) with no divisors but \(1\) and itself, and the
primes are the multiplicative atoms from which every integer is built. \B{PrimeQ} is the
test:

\begin{codepairs}
\pair{number-theory/primes/1}{\(97\) is prime.}
\pair{number-theory/primes/2}{\(91=7\cdot13\) is not---a number that looks prime until you
try to factor it.}
\pair{number-theory/primes/3}{\B{Prime} indexes the primes directly: the hundredth prime
is \(541\).}
\pair{number-theory/primes/4}{\B{PrimePi} counts them: there are \(25\) primes up to
\(100\).}
\end{codepairs}

The counting function \(\pi(x)\) is where number theory meets analysis. The Prime Number
Theorem\index{Prime Number Theorem} says \(\pi(x)\sim x/\ln x\), and Mathilda lets you
watch the approximation at \(x=10^{6}\):

\begin{codepairs}
\pair{number-theory/primes/5}{The exact count below a million.}
\pair{number-theory/primes/6}{The theorem's estimate \(x/\ln x\)---low by about eight
percent here, an error the theorem promises shrinks (slowly) to zero.}
\pair{number-theory/primes/7}{\B{NextPrime} finds the next prime above its argument,
instantly even past a million.}
\pair{number-theory/primes/8}{A negative second argument walks the other way, so this is
the largest prime \emph{below} \(100\)---there is no separate \texttt{PreviousPrime}.}
\end{codepairs}

\usagebox{PrimeQ}
\usagebox{PrimePi}

\B{PrimePi} is not one algorithm but a family, selectable through \texttt{Method}. Below a
threshold Mathilda simply sieves; above it, counting every prime is hopeless and it uses
a combinatorial method that never enumerates them---all agreeing, of course, on the
answer:

\begin{codepairs}
\pair{number-theory/primes/9}{Legendre's sieve-free recurrence.}
\pair{number-theory/primes/10}{The modern Del\'eglise--Rivat method---same count,
reachable to far larger \(x\).}
\end{codepairs}

A last twist: the notion of ``prime'' depends on the ring. In the Gaussian
integers\index{Gaussian integers} \(\mathbb{Z}[i]\) a rational prime may split, and
\B{PrimeQ} tests this too under \texttt{GaussianIntegers -> True}:

\begin{codepairs}
\pair{number-theory/primes/11}{\(5=(2+i)(2-i)\) is \emph{not} a Gaussian prime---it splits.}
\pair{number-theory/primes/12}{\(3\) stays prime in \(\mathbb{Z}[i]\); it is inert.}
\pair{number-theory/primes/13}{And \(2+3i\) is a Gaussian prime, its norm \(13\) being a
rational prime.}
\end{codepairs}

\calloutnote{Under the hood}{\B{PrimeQ} is not trial division. For machine-sized integers
it runs a deterministic Miller--Rabin test over a fixed witness set; for larger integers
it uses the Baillie--PSW test (a strong Fermat test composed with a Lucas test), for which
no counterexample is known below the ranges Mathilda will ever see. Both are
probabilistic in theory and, at these sizes, certain in practice~\cite{crandallpomerance}.}

Two smaller tools complete the kit. \mcode{RandomPrime} draws a prime uniformly at random
from a range---the first step of any key generation---and \mcode{PrimePowerQ} recognises
the numbers \(p^{k}\) that are ``prime up to a power'':

\begin{codepairs}
\pair{number-theory/primes-more/1}{Six primes below \(1000\). The draw is pseudorandom,
so the seed is fixed first with \B{SeedRandom}.}
\pair{number-theory/primes-more/2}{Reseeding with the same value reproduces the draw
exactly.}
\pair{number-theory/primes-more/3}{A \(100\)-digit prime, drawn and confirmed.}
\pair{number-theory/primes-more/4}{\(8=2^{3}\), \(9=3^{2}\) and \(-8\) are prime powers (signs
are ignored), \(12\) and \(1\) are not, and a rational counts when its reciprocal does.}
\pair{number-theory/primes-more/5}{A Gaussian argument is tested in \(\mathbb{Z}[i]\)
automatically: \((2+i)^{3}\) is a power of one Gaussian prime.}
\pair{number-theory/primes-more/6}{Asked in \(\mathbb{Z}[i]\), \(5=(2+i)(2-i)\) is \emph{not}
a prime power.}
\end{codepairs}

\usagebox{RandomPrime}
\usagebox{PrimePowerQ}

\calloutnote{Under the hood}{\mcode{RandomPrime} uses rejection sampling: it draws
uniformly from \(\{2\}\) together with the odd integers of the range and keeps the first
draw that passes the same primality test as \mcode{PrimeQ}. Every prime in the range is
therefore exactly equally likely---unlike the tempting shortcut of taking the
\mcode{NextPrime} after a random integer, which favours primes that follow long gaps. A range
that might contain no prime is checked for one first, so an empty range is reported with
a message rather than looped on forever. \mcode{PrimePowerQ} trial-divides, reduces the
number to its primitive root \(m^{1/k}\), and makes exactly one primality test on that
root.}

Counting primes is where the algorithms of this subsection earn their keep. The
\mcode{PrimePi} values below are far beyond any sieve that lists the primes one by one, yet
each takes well under a second:

\begin{codepairs}
\pair{number-theory/primes-more/7}{\(\pi(10^{12})\): there are \(37{,}607{,}912{,}018\) primes
below a trillion. \B{AbsoluteTiming} returns the elapsed time alongside the value.}
\pair{number-theory/primes-more/8}{And that count took less than a second.}
\pair{number-theory/primes-more/9}{The billionth prime, by way of \(\pi\).}
\end{codepairs}

\calloutnote{Performance}{The default \mcode{PrimePi} is a tuned
Del\'eglise--Rivat\index{Deleglise--Rivat method@Del\'eglise--Rivat method} algorithm, which
counts primes in roughly \(x^{2/3}\) time by a combinatorial sum over ``leaves'' instead of
by enumeration. On this book's build machine (an eight-core Intel i9-9880H) \(\pi(10^{12})\)
takes about \(0.2\)~s and \(\pi(10^{13})\) about \(1.2\)~s. \mcode{Prime}\([n]\) inverts the
count: it estimates the \(n\)th prime from Riemann's \(R(x)\approx\pi(x)\), makes one exact
\(\pi\) call at the estimate, and sieves only the short gap that remains, so
\(\mcode{Prime}[10^{9}]\) takes about \(0.02\)~s and \(\mcode{Prime}[10^{11}]\) about half a
second.}

\subsection{Integer factorization}
\label{ssec:nt-factor}

Testing primality is easy; \emph{factoring} a composite is the hard problem on which much
of modern cryptography rests. \B{FactorInteger} returns the prime factorization as a list
of \(\{\text{prime},\text{exponent}\}\) pairs:

\begin{codepairs}
\pair{number-theory/factorization/1}{\(360=2^{3}\cdot3^{2}\cdot5\).}
\pair{number-theory/factorization/2}{The fifth Fermat number \(2^{32}+1\). Fermat
conjectured all such numbers prime; Euler found this factor \(641\) in
\(1732\)\index{Fermat number}, and Mathilda recovers his result at once.}
\pair{number-theory/factorization/3}{A pandigital composite peeled apart into eight primes.}
\pair{number-theory/factorization/4}{A \(19\)-digit semiprime---two ten-digit primes---
factored in milliseconds. Its size is trivial for Mathilda but hints at the wall that RSA
hides behind: multiply the primes to hundreds of digits and no known method finishes.}
\end{codepairs}

\usagebox{FactorInteger}

The single call above is a whole cascade of algorithms, and \texttt{Method} lets you
name one explicitly:

\begin{codepairs}
\pair{number-theory/factorization/5}{The same semiprime, factored by Pollard's rho method
on demand\index{Pollard rho method}.}
\pair{number-theory/factorization/6}{\B{IntegerExponent} reads off the multiplicity of a
base directly: \(10^{6}\parallel10^{6}\).}
\pair{number-theory/factorization/7}{The power of \(2\) dividing \(360\).}
\pair{number-theory/factorization/8}{\B{SquareFreeQ}\index{squarefree}: \(30=2\cdot3\cdot5\) has no repeated
prime.}
\pair{number-theory/factorization/9}{\(12=2^{2}\cdot3\) does, so it is not square-free.}
\end{codepairs}

\calloutnote{Theory}{Factoring algorithms come in two kinds. \emph{Special-purpose}
methods find a factor \(p\) quickly when \(p\) has some property, whatever the size of
\(n\): Pollard's rho\index{Pollard rho method} takes about \(\sqrt p\) steps, so it is fast
for small \(p\); Pollard's \(p-1\)\index{Pollard p-1 method@Pollard \(p-1\) method} and
Williams' \(p+1\)\index{Williams p+1 method@Williams \(p+1\) method} succeed when
\(p-1\) or \(p+1\) has only small prime factors; Fermat's
method\index{Fermat's factoring method} when the two factors are close together; and
Lenstra's elliptic-curve method (ECM)\index{elliptic curve method}, which replaces the
fixed group \((\mathbb{Z}/p)^{\times}\) of \(p-1\) by the group of a random elliptic curve,
gets a fresh chance at a smooth group order with every curve, so its running time
depends on the size of \(p\) rather than of \(n\)~\cite{lenstraecm}. \emph{General-purpose}
methods ignore the factors' shape and pay by the size of \(n\): the continued-fraction
method (CFRAC) and the quadratic sieve\index{quadratic sieve} both collect many
congruences \(x^{2}\equiv y\pmod n\) with \(y\) smooth, then combine them by linear
algebra over \(\mathrm{GF}(2)\) into \(X^{2}\equiv Y^{2}\pmod n\), whence
\(\gcd(X-Y,n)\) is a factor about half the time. Their cost is subexponential in the
number of digits~\cite{crandallpomerance}.}

\mcode{FactorInteger} is not one algorithm but a family of them, and every member is also a
builtin of its own, living in the \mcode{FactorInteger\textasciigrave} context. Each
takes the method's tuning parameters as string-named options; the same call can be
spelled \mcode{Method -> \{"Name", "Sub" -> v\}} on \mcode{FactorInteger} itself. Choosing
the right tool for a number's \emph{shape} matters far more than its size:

\begin{codepairs}
\pair{number-theory/factor-methods/1}{\B{Names} lists the entry points: the classical
methods, three aliases (\mcode{QS}, \mcode{SIQS}, \mcode{SQUFOF}), and seven heads for the
rational base descent of \S\ref{ssec:nt-rbd}.}
\pair{number-theory/factor-methods/2}{The Fermat number \(F_{7}=2^{128}+1\) by ECM. Its
\(17\)-digit factor took Morrison and Brillhart a landmark computation in \(1970\).}
\pair{number-theory/factor-methods/3}{The same method through \texttt{Method}, here with
the stage-one smoothness bound \(B_{1}\) set by hand.}
\pair{number-theory/factor-methods/4}{Two \(31\)-digit primes a million apart: Fermat's
method writes \(n=a^{2}-b^{2}\) with \(a\) just above \(\sqrt n\), and finds this in its
first few steps.}
\pair{number-theory/factor-methods/5}{But on lopsided factors Fermat is hopeless. It
gives up when its budget runs out and returns the composite unfactored (the
\mcode{FactorInteger::\allowbreak nofac} message is silenced by \B{Quiet}).}
\pair{number-theory/factor-methods/6}{The same number by \(p-1\): the factor \(2^{61}-1\)
has \(p-1=2\cdot3^{2}\cdot5^{2}\cdot7\cdot11\cdot13\cdot31\cdot41\cdot61\cdot151\cdot331\cdot1321\),
entirely made of small primes, so it falls out at once.}
\pair{number-theory/factor-methods/7}{A balanced \(60\)-digit semiprime has no shape to
exploit; this is the quadratic sieve's territory, about half a second here.}
\pair{number-theory/factor-methods/8}{The eighth Fermat number, \(2^{256}+1\), whose
factorization by Brent and Pollard in \(1980\) was a celebrated success for the rho
method. \texttt{Automatic} finds the \(16\)-digit factor in a fraction of a second.}
\end{codepairs}

Two further forms of the call round out the interface. A second argument \(k\) asks for a
\emph{partial} factorization with at most \(k\) entries, which is useful when you need
only the small primes of a number whose large cofactor is expensive. The entries for the
largest primes are kept and the smallest are merged into one composite entry. And
\texttt{GaussianIntegers -> True} factors over \(\mathbb{Z}[i]\), as a Gaussian argument
does automatically:

\begin{codepairs}
\pair{number-theory/factor-forms/1}{\(10!=2^{8}\,3^{4}\,5^{2}\,7\) in full.}
\pair{number-theory/factor-forms/2}{At most three entries: \(5^{2}\) and \(7\) are kept,
and \(2^{8}3^{4}=20736\) is merged into one.}
\pair{number-theory/factor-forms/3}{A Gaussian integer factors over \(\mathbb{Z}[i]\) on its
own: \(3+4i=(2+i)^{2}\).}
\pair{number-theory/factor-forms/4}{\(5=-i\,(1+2i)(2+i)\). The unit comes first, then the
primes in the first quadrant---the splitting of \S\ref{ssec:nt-primes}, made explicit.}
\end{codepairs}

\calloutnote{Under the hood}{With \texttt{Method -> Automatic}, \mcode{FactorInteger} runs a
cascade of tiers, cheapest first. Below \(2^{64}\) a separate word-size engine does
everything in machine arithmetic, with no big-integer allocation: Montgomery
multiplication, a deterministic Miller--Rabin test, Pollard rho, and a ``micro-ECM'' on
Suyama curves, with Shanks' square forms (SQUFOF) as a backstop. Above \(2^{64}\) the
tiers are trial division to \(2^{16}\), a perfect-power test, Pollard rho (Brent's
variant), a light two-word ECM, a Fermat probe for near-square factors, one
\(p-1\) pass, one \(p+1\) pass, and then ECM at increasing factor sizes of \(15\), \(20\),
\(25\), \dots\ digits (through GMP-ECM). After every split both parts re-enter at the tier
that succeeded, so a number rich in small factors keeps being peeled by the cheap
tiers. Once the search effort spent---measured as the size of factor the earlier tiers
would probably have found---reaches \(2/7\) of the composite's digit count, the point past
which sieving is cheaper than searching on, a composite of up to \(100\)
digits is handed to Mathilda's own self-initialising quadratic
sieve\index{quadratic sieve!self-initialising}: a Knuth--Schroeppel multiplier, Gray-code switching between
polynomials, a cache-sized block sieve, single large primes and, from \(72\) digits,
double large primes combined by finding cycles in a graph, with block
Lanczos\index{block Lanczos algorithm} for the \(\mathrm{GF}(2)\) linear algebra. The
\mcode{Lehman}, \mcode{CFRAC} and \mcode{Dixon} methods are not part of the cascade; they
are there to be called by name.}

\calloutnote{Under the hood}{Every random choice in \mcode{FactorInteger}---rho constants, ECM
curves, \(p+1\) seeds, sieve polynomials---comes from a generator seeded by \(n\)
itself\index{factorization!deterministic seeding}, so the answer, the time it takes, and
whether a hard number is left unfactored depend on \(n\) alone. The quadratic sieve from
\(45\) digits, ECM above \(2^{160}\), and \(p+1\), CFRAC, Dixon and the rational base
descent all spread their work across every core (up to sixteen threads), but the results
of the workers are replayed in a fixed order, so the answer does not depend on the thread
count either. \B{TimeConstrained} budgets are measured in elapsed time, and a parallel
region checks the deadline often enough that an abort is taken promptly.}

How fast is it? The table below is drawn from the shootout in the repository's
\texttt{benchmarks/\allowbreak 98-factorinteger-shootout/\allowbreak REPORT.md}, run at version \(0.453\) on the
same eight-core i9-9880H that built this book's transcripts. PARI/GP~2.17 and FLINT~3.6
ran single-threaded; Mathilda used its threads, and ``---'' marks a run that timed out
(600~s for PARI and FLINT; the reference-system column is from an earlier run with a
120~s limit). Batch rows give the total for the whole group.

\begin{center}
\small
\begin{tabular}{@{}lrrrr@{}}
\toprule
Input & Mathilda & PARI/GP & FLINT & reference \\
\midrule
\(500\) semiprimes of \(63\) bits & \textbf{38 ms} & 154 ms & 126 ms & 585 ms \\
\(20\) semiprimes of \(96\) bits & 77 ms & \textbf{40 ms} & 127 ms & 244 ms \\
\(30\)-digit semiprime & 6.4 ms & \textbf{3.0 ms} & 6.2 ms & 62 ms \\
\(50\)-digit semiprime & \textbf{76 ms} & 279 ms & 358 ms & 6.2 s \\
\(60\)-digit semiprime & \textbf{0.61 s} & 3.2 s & 6.5 s & --- \\
\(70\)-digit semiprime & \textbf{3.1 s} & 27 s & 27 s & --- \\
\(80\)-digit semiprime & \textbf{32 s} & 494 s & 346 s & --- \\
\(90\)-digit semiprime & \textbf{288 s} & --- & --- & --- \\
\(25\)-digit factor by ECM & \textbf{3.2 s} & 188 s & 207 s & --- \\
\(F_{8}=2^{256}+1\) & \textbf{0.19 s} & 0.84 s & 0.34 s & 0.98 s \\
\bottomrule
\end{tabular}
\end{center}

\calloutnote{Performance}{Two honest qualifications to that table. First, the threads:
single-threaded, Mathilda's sieve still factors the \(60\)-digit case in about
\(0.4\)~s against PARI's \(3.4\)~s, and the \(80\)-digit case in \(145\)~s against
\(320\)--\(490\)~s, so the lead is algorithmic, not just parallel. Second, the rows it
loses are the small ones, where fixed costs dominate: at \(30\) digits the cascade spends
milliseconds on rho and ECM before it switches to the sieve, and on \(96\)-bit
semiprimes a minority of inputs slip past the light ECM tier. Above all, the asymmetry
between \S\ref{ssec:nt-primes} and this section is the whole point. Deciding that a
\(300\)-digit number is composite takes milliseconds; producing its factors is, as far as
anyone knows, infeasible. RSA (\S\ref{ssec:nt-rsa}) is exactly this gap, weaponised.}

\subsection{Blake's rational base descent}
\label{ssec:nt-rbd}

The general-purpose methods are blind to structure, and the special-purpose ones each
look for one particular kind. \emph{Rational base descent}\index{rational base descent}
(RBD) is a special-purpose method of this second kind, devised by Sam Blake, the author of
Mathilda~\cite{blake-rbd,blake-lattice-rbd}. Its target is a number \(n=PM\) whose
\emph{larger} factor lies close to a power of a rational number,
\[
  P \approx c\left(\frac{a}{b}\right)^{k},\qquad a>b\ge1,\ \gcd(a,b)=1,
\]
with a modest coefficient \(c\); the smaller factor \(M\) can be anything. Primes of this
shape are not something an honest key generator produces, but they are exactly what a
careless or malicious one might, and the method also explains why integers with
algebraic structure, such as \(2^{n}\pm1\), split along their algebraic factors.

The idea is a descent. Start from \(Q=n\) and repeatedly divide by the base, discarding
fractions:
\[
  Q \;\leftarrow\; \left\lfloor \frac{b\,Q}{a} \right\rfloor .
\]
After \(k\) steps the factor \(P\approx c(a/b)^{k}\) has been divided down to about \(c\),
so \(Q\) has become approximately \(cM\)---an unknown multiple of the hidden factor. If it
is close enough, then for some small offset \(j\) the gcd \(\gcd(Q-j,\,n)\) equals \(M\)
(or \(M\gcd(c,P)\)), and the number is split. One descent passes through every depth, so
it tests every exponent \(k\) and every coefficient \(c\) at once, in time polynomial in
\(\log n\).

\begin{codepairs}
\pair{number-theory/rbd/1}{A planted instance: \(P\) is the first prime after
\(3\,(22/7)^{150}\), a \(76\)-digit number.}
\pair{number-theory/rbd/2}{The second factor \(M\) is an arbitrary \(30\)-digit
prime, the first after \(7\cdot10^{29}\).}
\pair{number-theory/rbd/3}{The defect \(\Delta=P-\lfloor c(a/b)^{k}\rfloor\) is only \(46\).}
\pair{number-theory/rbd/4}{The product has \(105\) digits.}
\pair{number-theory/rbd/5}{Told the base \(22/7\), the paper's reference descent
(\mcode{Classic}) splits it instantly.}
\pair{number-theory/rbd/6}{The default variant does not even need the base: it enumerates
candidate bases and finds a working one in about ten milliseconds.}
\pair{number-theory/rbd/7}{\texttt{Automatic} \mcode{FactorInteger} does not use RBD. Facing a
\(30\)-digit factor of a \(105\)-digit number---too large for its sieve, a hard case
for ECM---it is still searching when two seconds run out.}
\end{codepairs}

\calloutnote{Theory}{Why the descent lands near \(cM\): write each step as \(Q_{i}=(b/a)Q_{i-1}-\epsilon_{i}\)
with \(\epsilon_{i}\in[0,1)\). Unrolling, \(Q_{k}=(b/a)^{k}n-E_{k}\) with
\(0\le E_{k}<a/(a-b)\), a geometric series. Substituting \(n=PM\) and
\(P=\lfloor c(a/b)^{k}\rfloor+\Delta\) gives
\[
  Q_{k}-cM \;=\; \Delta\left(\tfrac{b}{a}\right)^{k}M \;-\; \delta\left(\tfrac{b}{a}\right)^{k}M \;-\; E_{k},
  \qquad \delta=\{c(a/b)^{k}\}\in[0,1).
\]
Blake's theorem~\cite{blake-rbd} follows by bounding the three terms: if \((a/b)^{k}>M\)
and \(|\Delta|<(a/b)^{k}/M\), then \(|Q_{k}-cM|\le J=\lceil2+a/(a-b)\rceil\), so a window of
\(2J+1\) gcds at each depth cannot miss. The same bound shows the structure must sit on
the \emph{larger} factor, and that against a balanced RSA modulus the tolerance
\((a/b)^{k}/M\) shrinks to about one, so an unstructured \(n\) needs on the order of
\(\sqrt n\) candidate bases---no threat to RSA. The bases that work form a ``comb'' of
isolated teeth\index{rational base descent!comb of bases} of relative width about
\(1/(kM)\): nothing signals a nearby base, so bases are enumerated, breadth-first over the
Stern--Brocot tree\index{Stern--Brocot tree}, which lists every positive rational once,
coarse to fine.}

The theorem's window is the method's weakness. For a balanced semiprime the tolerance is
about one: if \(P\) misses \(c(a/b)^{k}\) by more than a few units, the descent's \(Q\) lands
too far from \(cM\) and no gcd sees it. The companion paper~\cite{blake-lattice-rbd}
replaces the gcd window by a \emph{lattice} detector: Howgrave-Graham's method for
approximate common divisors\index{approximate common divisor}, which uses LLL lattice
reduction\index{LLL lattice reduction} to recover \(cM\) from any \(Q\) within a distance
\(X\approx n^{\beta^{2}}\) when \(M\ge n^{\beta}\). The next session plants a miss of
almost \(2^{28}\) in a balanced \(60\)-digit semiprime:

\begin{codepairs}
\pair{number-theory/rbd-lattice/1}{\(P\) is now the first prime after
\(3\,(22/7)^{59}+2^{28}\)---deliberately far from the power.}
\pair{number-theory/rbd-lattice/2}{\(M\) is a prime of about the same size.}
\pair{number-theory/rbd-lattice/3}{The \(60\)-digit product.}
\pair{number-theory/rbd-lattice/4}{The miss \(Q_{59}-cM\) of the exact descent: about
\(2.4\times10^{8}\), against a window of a few units.}
\pair{number-theory/rbd-lattice/5}{So the classic descent fails and returns \(n\)
unfactored.}
\pair{number-theory/rbd-lattice/6}{The lattice variant runs one LLL reduction at each depth
of the band where \(cM\) can lie, and finds it.}
\pair{number-theory/rbd-lattice/7}{Knowing \(a/b\), \(c\) and \(k\) as well turns the search
into a single detector call at \(\lfloor c(a/b)^{k}\rfloor\): this is Coppersmith-style
factoring with the high bits of \(P\) known\index{Coppersmith's method}, here spelled with
\texttt{Method} and a \texttt{"Variant"} suboption.}
\end{codepairs}

There are six variants, each selected by \texttt{"Variant" -> name} and each also a head
of its own, \mcode{FactorInteger\textasciigrave BlakeRationalBaseDescent}\emph{Name}:

\begin{itemize}
\item \texttt{"Exact"} (the default) runs the descent on \(Sn\) for a power of two
\(S\ge2a/(a-b)\), so the target always lies among four consecutive values whatever the
base; it skips depths above \(\sqrt{c_{\max}n}\), where \(cM\) cannot lie, and batches
its gcds into one per block of steps.
\item \texttt{"Classic"} is the paper's reference algorithm, a gcd window of \(2J+1\)
values at every depth.
\item \texttt{"AllHits"} runs the descent to the end and refines every non-trivial gcd it
meets into pairwise coprime divisors---the mode for Cunningham numbers \(b^{n}\pm1\) and
repunits, whose base-\(b\) truncations inherit their algebraic factors.
\item \texttt{"Lattice"} swaps the gcd window for the Howgrave-Graham detector.
\item \texttt{"TwoPass"}, the papers' practical algorithm, is \texttt{"Exact"} over the
enumerated bases followed by \texttt{"Lattice"} over a coarse grid of bases.
\item \texttt{"Coppersmith"} takes a known base, coefficient and exponent and makes one
detector call.
\end{itemize}

\calloutnote{Under the hood}{The descent core is a small, reentrant library over GMP, with
FLINT's LLL when it is available and an exact native LLL otherwise. A many-base search
runs blocks of \(2048\) bases on every thread, and a lattice base runs its depths two per
thread; in both, the lowest-indexed success wins, so the answer does not depend on the
thread count. Suboptions set the search: \texttt{"Base"} (one base or a list),
\texttt{"Order"} (Stern--Brocot or by \(a+b\)), \texttt{"Height"} and \texttt{"Bases"}
(the enumeration limits), \texttt{"MaxCoefficient"}, \texttt{"Multiplier"}, and the
lattice parameters \texttt{"Tolerance"}, \texttt{"Beta"}, \texttt{"Dimension"} and
\texttt{"Multiplicity"}. The method's \texttt{?} documentation, which is unusually long,
explains each one. The repository's tests reproduce every example in the two papers; the
papers' \(2048\)-bit challenge modulus is factored from scratch, with no base given, in
about two minutes on sixteen threads.}

\subsection{Arithmetic functions}
\label{ssec:nt-arith}

Once an integer is factored, a whole family of \emph{arithmetic
functions}\index{arithmetic function}---functions of a single integer, defined through its
factorization---becomes computable~\cite{hardywright}. Most are \emph{multiplicative}\index{multiplicative
function}: their value on \(mn\) is the product of their values on \(m\) and \(n\) when
\(m,n\) are coprime. The most basic list or count the divisors:

\begin{codepairs}
\pair{number-theory/arithmetic-functions/1}{\B{Divisors} lists every positive divisor of
\(36\), in order.}
\pair{number-theory/arithmetic-functions/2}{\B{DivisorSigma} with first argument \(0\)
counts them: \(36\) has nine divisors.}
\pair{number-theory/arithmetic-functions/3}{With first argument \(1\) it sums them;
\(\sigma_1(28)=56\), a fact we return to below.}
\end{codepairs}

Euler's totient \(\varphi(n)\) counts the integers in \(1,\dots,n\) coprime to \(n\)---
equivalently, the size of the group of units modulo \(n\):

\begin{codepairs}
\pair{number-theory/arithmetic-functions/4}{\(\varphi(36)=12\).}
\pair{number-theory/arithmetic-functions/5}{The same number from the product formula
\(n\prod_{p\mid n}(1-1/p)\)---this is what ``multiplicative'' buys you.}
\end{codepairs}

The remaining functions read off the shape of the factorization. The M\"obius function
\(\mu(n)\)\index{Mobius function@M\"obius function} is \(0\) when a square divides \(n\)
and \(\pm1\) otherwise; Liouville's \(\lambda(n)\), \(\Omega(n)\), and \(\nu(n)\) count
prime factors with and without multiplicity:

\begin{codepairs}
\pair{number-theory/arithmetic-functions/6}{\(\mu(30)=(-1)^{3}\): three distinct primes,
none repeated.}
\pair{number-theory/arithmetic-functions/7}{\(\mu(12)=0\): the square \(4\) divides it.}
\pair{number-theory/arithmetic-functions/8}{\(\lambda(12)=(-1)^{\Omega(12)}=(-1)^{3}=-1\).}
\pair{number-theory/arithmetic-functions/9}{\B{PrimeOmega} counts primes \emph{with}
multiplicity: \(360=2^3\cdot3^2\cdot5\) has \(\Omega=6\).}
\pair{number-theory/arithmetic-functions/10}{\B{PrimeNu} counts \emph{distinct} primes:
\(\nu(360)=3\).}
\pair{number-theory/arithmetic-functions/11}{Like the divisor functions, these extend to
\(\mathbb{Z}[i]\): the Gaussian divisors of \(5\) exhibit its splitting \(5=(2+i)(2-i)\).}
\end{codepairs}

\usagebox{DivisorSigma}
\usagebox{EulerPhi}
\usagebox{MoebiusMu}

These functions are not curiosities; the divisor sum is the gateway to one of the oldest
questions in mathematics. A \emph{perfect number}\index{perfect number} equals the sum of
its proper divisors---equivalently, \(\sigma_1(n)=2n\)---and Mathilda finds the first four
by brute force in an instant:

\begin{codepairs}
\pair{number-theory/perfect/1}{Every perfect number below \(10^{4}\): \(6,28,496,8128\).}
\pair{number-theory/perfect/2}{Euclid's formula \(2^{p-1}(2^{p}-1)\) for \(p=2,3,5,7\)
reproduces exactly that list\index{Euclid--Euler theorem}.}
\end{codepairs}

The match is the Euclid--Euler theorem: \(2^{p-1}(2^{p}-1)\) is perfect precisely when
\(2^{p}-1\) is a \emph{Mersenne prime}\index{Mersenne prime}. So the even perfect numbers
are in exact correspondence with the Mersenne primes, and testing one is testing the
other:

\begin{codepairs}
\pair{number-theory/perfect/3}{\(2^{13}-1=8191\) is a Mersenne prime, giving the fifth
perfect number.}
\pair{number-theory/perfect/4}{But \(2^{11}-1\) is \emph{not} prime---the exponent being
prime is necessary, not sufficient.}
\pair{number-theory/perfect/5}{Indeed \(2^{11}-1=2047=23\cdot89\).}
\end{codepairs}

\subsection{Modular arithmetic and congruences}
\label{ssec:nt-modular}

To work modulo \(n\) is to care only about remainders. \mcode{Mod} (met in
\S\ref{sec:arithmetic}, and always returning a residue with the sign of the modulus) is
the entry point; the number-theoretic weight sits in what you can do inside a modulus. The
first essential is fast modular exponentiation, because forming a huge power and
\emph{then} reducing it is wasteful when you can reduce at every step:

\begin{codepairs}
\pair{number-theory/modular/1}{Forming \(7^{100}\) in full and reducing---correct, but it
builds an \(85\)-digit integer first.}
\pair{number-theory/modular/2}{\B{PowerMod} gives the same residue without ever forming
the giant power.}
\pair{number-theory/modular/3}{A far larger exponent, still immediate.}
\end{codepairs}

\B{PowerMod} carries two overloads that pack a surprising amount of number theory. A
negative exponent asks for the modular \emph{inverse}, and a fractional one for a modular
\emph{root}:

\begin{codepairs}
\pair{number-theory/modular/4}{The inverse of \(3\) modulo \(7\) is \(5\), since
\(3\cdot5=15\equiv1\).}
\pair{number-theory/modular/5}{A modular square root: \(3^{2}=9\equiv2\pmod 7\), computed
by the Tonelli--Shanks algorithm\index{Tonelli--Shanks algorithm}.}
\end{codepairs}

The \emph{multiplicative order} of \(a\) modulo \(n\) is the least power returning \(a\) to
\(1\); when it reaches the full \(\varphi(n)\), \(a\) is a \emph{primitive root}\index{primitive
root}, a generator of the units:

\begin{codepairs}
\pair{number-theory/modular/6}{\(2\) has order \(3\) modulo \(7\) (\(2^{3}=8\equiv1\)).}
\pair{number-theory/modular/7}{\(3\) has order \(6\)---the full group---so it is a
primitive root.}
\pair{number-theory/modular/8}{\B{PrimitiveRoot} returns the least one, \(3\).}
\pair{number-theory/modular/9}{\B{PrimitiveRootList} gives all of them; there are
\(\varphi(\varphi(7))=\varphi(6)=2\).}
\pair{number-theory/modular/10}{Euler's theorem in one line: \(3^{\varphi(10)}\equiv1\),
because \(3\) is coprime to \(10\).}
\end{codepairs}

The order also frames the hardest question in this circle of ideas. Reversing it---given a
base \(a\) and a target \(b\), find the exponent \(x\) with \(a^{x}\equiv b\pmod n\)---is the
\emph{discrete logarithm problem}\index{discrete logarithm problem}, and the three-argument
form of \B{MultiplicativeOrder} solves exactly that, returning the least exponent whose power
lands in a given residue set:

\begin{codepairs}
\pair{number-theory/modular/13}{The discrete logarithm of \(5\) to base \(3\) modulo \(7\):
the least \(x\) with \(3^{x}\equiv5\), namely \(x=5\) (indeed \(3^{5}=243\equiv5\)).}
\end{codepairs}

\calloutnote{Theory}{Computing the order of \(a\) is easy---it divides \(\varphi(n)\) and is
found by stripping prime factors from it. The discrete logarithm is the \emph{reverse}
question, and it is believed \emph{hard}: no polynomial-time algorithm is known over a
general prime field, and that presumed hardness is the foundation of Diffie--Hellman key
exchange and the ElGamal and DSA signature schemes\index{discrete logarithm problem!in
cryptography}. The best general methods---baby-step giant-step, Pollard's rho for logarithms,
Pohlig--Hellman (which reduces the problem to the prime-power factors of the group order),
and index calculus---are surveyed in~\cite{hac}. Mathilda's three-argument form simply walks
the orbit \(a^{1},a^{2},\dots\), so it is a teaching tool for small moduli, not a break of
the problem at cryptographic sizes.}

Finally, the \B{JacobiSymbol}\index{Jacobi symbol} \(\left(\tfrac{a}{n}\right)\) encodes quadratic
residues\index{quadratic reciprocity}---whether \(a\) is a square modulo \(n\):

\begin{codepairs}
\pair{number-theory/modular/11}{\(\left(\tfrac{2}{7}\right)=+1\): \(2\) is a square mod
\(7\) (indeed \(3^{2}\equiv2\), as above).}
\pair{number-theory/modular/12}{\(\left(\tfrac{3}{7}\right)=-1\): \(3\) is a non-residue.}
\end{codepairs}

\usagebox{PowerMod}
\usagebox{MultiplicativeOrder}
\usagebox{PrimitiveRoot}
\usagebox{JacobiSymbol}

\calloutnote{Theory}{Two theorems govern this section. Fermat's little
theorem\index{Fermat's little theorem}---\(a^{p-1}\equiv1\pmod p\) for a prime \(p\nmid
a\)---is the special case \(n=p\) of \emph{Euler's theorem}\index{Euler's theorem},
\(a^{\varphi(n)}\equiv1\pmod n\) whenever \(\gcd(a,n)=1\). Together with the fact that the
units modulo a prime form a \emph{cyclic} group---generated by a primitive root---they are
the whole basis of what follows, and of RSA in \S\ref{ssec:nt-rsa}~\cite{irelandrosen}.}

The \B{ChineseRemainder} builtin solves a system of simultaneous
congruences\index{Chinese Remainder Theorem}\index{congruence!system of} in a
single step. Given residues \(r_i\) and moduli \(n_i\), it returns the smallest
non-negative \(x\) with \(x\equiv r_i\pmod{n_i}\) for every \(i\)---unique modulo
the moduli's least common multiple:

\begin{codepairs}
\pair{number-theory/crt/1}{The three moduli are listed.}
\pair{number-theory/crt/2}{The matching residues are listed too.}
\pair{number-theory/crt/3}{The unique \(x\) in \([0,105)\) that is \(\equiv2,3,2\)
modulo \(3,5,7\)---the residue \(23\).}
\pair{number-theory/crt/4}{And it checks: \(23\) leaves exactly those remainders.}
\end{codepairs}

It is instructive to see the classical recipe the builtin implements---the
Chinese Remainder Theorem\index{Chinese Remainder Theorem!construction} is a
short composition of the tools above. With the combined modulus \(N=\prod n_i\),
each cofactor \(N/n_i\) is invertible modulo its own \(n_i\) (the \(n_i\) being
coprime), and the \B{PowerMod} inverses assemble a weighted sum that reproduces
every residue at once---the recipe of Sun Zi\index{Sun Zi}:

\begin{codepairs}
\pair{number-theory/crt/5}{The combined modulus \(3\cdot5\cdot7=105\).}
\pair{number-theory/crt/6}{Each cofactor \(N/n_i\).}
\pair{number-theory/crt/7}{Each cofactor's inverse modulo its own \(n_i\).}
\pair{number-theory/crt/8}{The weighted sum, reduced---the same residue \(23\).}
\end{codepairs}

Two conveniences round out the builtin. A third argument shifts the search
window, returning the smallest solution at least that large; and the moduli need
not be coprime---\B{ChineseRemainder} returns unevaluated when the congruences
are inconsistent (here \(\gcd(6,10)=2\), yet \(1\not\equiv2\pmod2\)):

\begin{codepairs}
\pair{number-theory/crt/9}{The smallest solution \(\ge100\): namely \(23+105=128\).}
\pair{number-theory/crt/10}{No \(x\) satisfies both congruences, so the call
stands unevaluated.}
\end{codepairs}

\usagebox{ChineseRemainder}

\subsection{Continued fractions}
\label{ssec:nt-contfrac}

A continued fraction\index{continued fraction} writes a number as
\(a_0+\cfrac{1}{a_1+\cfrac{1}{a_2+\cdots}}\), a form that captures rational approximation
better than any decimal. \B{ContinuedFraction} produces the list of partial quotients and
\B{FromContinuedFraction} folds it back:

\begin{codepairs}
\pair{number-theory/continued-fractions/1}{\(415/93=[4;2,6,7]\). These quotients are
exactly Euclid's from \S\ref{ssec:nt-euclid}, run on \(415\) and \(93\).}
\pair{number-theory/continued-fractions/2}{Folding the list back recovers the rational
exactly---the two operations are inverse.}
\end{codepairs}

Irrationals have \emph{infinite} expansions, so you ask for a fixed number of terms---and
the classical patterns appear at once:

\begin{codepairs}
\pair{number-theory/continued-fractions/3}{\(\sqrt2=[1;\overline{2}]\), purely periodic.}
\pair{number-theory/continued-fractions/4}{\(\sqrt7\) has period \(1,1,1,4\).}
\pair{number-theory/continued-fractions/5}{Asked without a length, Mathilda returns the
period in a nested list---every quadratic irrational is \emph{eventually periodic}, which
is Lagrange's theorem\index{quadratic irrational}.}
\pair{number-theory/continued-fractions/6}{\(\pi\) shows no pattern, but its early terms
are the key to good approximations.}
\end{codepairs}

Truncating \(\pi\)'s expansion just after the large term \(15\) and folding it back gives
the approximation known since antiquity:

\begin{codepairs}
\pair{number-theory/continued-fractions/7}{\(355/113\), from only four terms.}
\pair{number-theory/continued-fractions/8}{It agrees with \(\pi\) to better than
\(3\times10^{-7}\)---the best rational approximation with any denominator under a thousand.}
\end{codepairs}

\usagebox{ContinuedFraction}
\usagebox{FromContinuedFraction}

\calloutnote{Theory}{A large partial quotient signals an unusually good rational
approximation just before it: truncating there makes the discarded tail
\(1/(a_{k+1}+\cdots)\) tiny. \(\pi\)'s term \(292\) is why \(355/113\) is so
extraordinarily accurate. The convergents obtained by truncation are, in a precise sense,
the \emph{best} rational approximations to a number\index{continued fraction!convergent}~\cite{khinchin}---the theory that also governs Pell's
equation, which \S\ref{ssec:alg-diophantine} solves through the continued fraction of
\(\sqrt D\).}

\subsection{Digits and bases}
\label{ssec:nt-digits}

The written form of an integer---its digits in a given base---is itself a number-theoretic
object, and Mathilda exposes it directly. \B{IntegerDigits} and \B{FromDigits} convert both
ways, in any base:

\begin{codepairs}
\pair{number-theory/digits/1}{The base-ten digits of \(123456\).}
\pair{number-theory/digits/2}{In base two, \(255\) is eight ones.}
\pair{number-theory/digits/3}{In base sixteen it is two ``fifteens''.}
\pair{number-theory/digits/4}{\B{FromDigits} rebuilds the integer from its base-two digits.}
\pair{number-theory/digits/5}{\B{IntegerString} renders the digits as text, here the
familiar hexadecimal \texttt{ff}.}
\pair{number-theory/digits/6}{\B{IntegerLength} counts the digits without listing them.}
\end{codepairs}

The digit sum is the basis of the schoolroom rule of \emph{casting out nines}---a number
is congruent to its digit sum modulo \(9\), because \(10\equiv1\):

\begin{codepairs}
\pair{number-theory/digits/7}{\B{DigitSum} adds the base-ten digits: \(1+2+3+4+5+6=21\).}
\pair{number-theory/digits/8}{Its residue modulo \(9\) is taken.}
\pair{number-theory/digits/9}{It equals the number's own residue modulo \(9\). That is
casting out nines.}
\pair{number-theory/digits/10}{\B{DigitCount} tallies how often each digit occurs---one
\(1\), two \(2\)s, three \(3\)s, four \(4\)s (the final slot counts the zeros).}
\end{codepairs}

Base two earns its own function. \B{BitLength} returns the number of binary digits---the
position of the highest set bit---which is \B{IntegerLength} in base \(2\), and the natural
measure of how many bits a number occupies:

\begin{codepairs}
\pair{number-theory/digits/11}{\B{BitLength} of \(255\) is \(8\): the eight ones of pair~2
above.}
\pair{number-theory/digits/12}{It is exact for arbitrarily large integers---\(2^{100}\)
needs \(101\) bits, the leading one plus a hundred zeros.}
\end{codepairs}

Not every positional-looking notation is positional. \B{FromDigits} also reads Roman
numerals\index{Roman numerals}, an additive system with subtractive pairs (\texttt{IV},
\texttt{XC}) rather than a base:

\begin{codepairs}
\pair{number-theory/digits/13}{\(\mathrm{MMXXVI}=1000+1000+10+10+5+1=2026\).}
\end{codepairs}

\usagebox{IntegerDigits}

\subsection{Combinatorial number theory}
\label{ssec:nt-combinatorial}

Number theory shades into combinatorics through the factorials, binomial coefficients, and
partition functions---all computed exactly, all growing explosively. The factorial family
comes first:

\begin{codepairs}
\pair{number-theory/combinatorial/1}{\(20!\), exact to the last of its nineteen digits.}
\pair{number-theory/combinatorial/2}{The double factorial \(9!!=9\cdot7\cdot5\cdot3\cdot1\).}
\pair{number-theory/combinatorial/3}{The falling factorial \(10\cdot9\cdot8\), via
\B{FactorialPower}.}
\pair{number-theory/combinatorial/4}{\B{Binomial}: the number of \(3\)-subsets of a
\(10\)-set.}
\pair{number-theory/combinatorial/5}{Reduced modulo \(7\) it is \(1\)---as Lucas's theorem
predicts from the base-\(7\) digits of \(10\) and \(3\)\index{Lucas's theorem}.}
\end{codepairs}

The partition functions count the ways to write \(n\) as a sum of positive integers.
\B{IntegerPartitions} lists them and \B{PartitionsP} counts them, and the count grows so
fast that its exact evaluation is a small miracle of analysis:

\begin{codepairs}
\pair{number-theory/combinatorial/6}{The seven partitions of \(5\).}
\pair{number-theory/combinatorial/7}{Their count, \(p(10)=42\).}
\pair{number-theory/combinatorial/8}{\(p(100)=190569292\)---already far past listing.}
\pair{number-theory/combinatorial/9}{\(p(1000)\), a \(32\)-digit integer computed exactly.}
\pair{number-theory/combinatorial/10}{\B{PartitionsQ} counts partitions into
\emph{distinct} parts.}
\pair{number-theory/combinatorial/15}{A length specification restricts the parts: there
are \(7153\) ways to write \(100\) as a sum of exactly four positive integers.}
\end{codepairs}

\noindent That last count is answered by listing, but not by building every partition and
filtering: \mcode{IntegerPartitions} walks the restricted partitions directly, so the
list above costs about two milliseconds.

Finally the Fibonacci and Lucas numbers, the archetypal integer sequences:

\begin{codepairs}
\pair{number-theory/combinatorial/11}{\(F_{50}\).}
\pair{number-theory/combinatorial/12}{\(F_{200}\), a \(42\)-digit number, still instant.}
\pair{number-theory/combinatorial/13}{The tenth Lucas number.}
\pair{number-theory/combinatorial/14}{A jewel of an identity: \(\gcd(F_m,F_n)=F_{\gcd(m,n)}\),
here \(\gcd(F_{12},F_{18})=F_{6}\).}
\end{codepairs}

\usagebox{Binomial}
\usagebox{PartitionsP}
\usagebox{Fibonacci}

\calloutnote{Under the hood}{\B{PartitionsP} does not count partitions one by one---that
would be hopeless at \(n=1000\). It evaluates the Hardy--Ramanujan--Rademacher
series\index{Hardy--Ramanujan formula}~\cite{andrews}, a convergent sum of the analytic terms that gives
\(p(n)\) exactly after enough of them, rounded to the nearest integer. \B{Fibonacci} uses
fast doubling---computing \(F_{2k}\) and \(F_{2k+1}\) from \(F_k\) and \(F_{k+1}\)---so
\(F_n\) costs \(O(\log n)\) big-integer multiplications rather than \(n\) additions. The
close relatives \B{BernoulliB}, \B{EulerE}, and \B{Zeta}, whose values encode still deeper
number theory, belong to the special functions of \S\ref{sec:specialfns}.}

\subsection{Putting it together: textbook RSA}
\label{ssec:nt-rsa}

Almost every tool of this section meets in one place: the RSA cryptosystem\index{RSA},
whose security is exactly the factoring gap of \S\ref{ssec:nt-factor}. We use the historic
textbook primes \(p=61\), \(q=53\) so every step is legible; the session is continuous,
each line building on the last. First the public modulus and its secret totient:

\begin{codepairs}
\pair{number-theory/rsa/1}{Pick the first prime, \(p=61\).}
\pair{number-theory/rsa/2}{Pick the second prime, \(q=53\).}
\pair{number-theory/rsa/3}{The modulus \(n=pq=3233\) is public; its factors are the secret.}
\pair{number-theory/rsa/4}{The totient is \(\varphi(n)=3120\).}
\pair{number-theory/rsa/5}{For a product of two primes this is \((p-1)(q-1)\), as
knowing the factorization lets us confirm.}
\end{codepairs}

Next a public exponent coprime to \(\varphi(n)\), and the private exponent that inverts it:

\begin{codepairs}
\pair{number-theory/rsa/6}{The conventional public exponent is \(e=17\).}
\pair{number-theory/rsa/7}{It is coprime to \(\varphi(n)\), so it is invertible.}
\pair{number-theory/rsa/8}{\B{ExtendedGCD} already exposes the inverse in its B\'ezout
coefficient: \((-367)\cdot17+2\cdot3120=1\).}
\pair{number-theory/rsa/9}{The private key \(d\) is that inverse as a positive residue,
\(2753\) (indeed \(-367+3120\)).}
\pair{number-theory/rsa/10}{Confirming the defining relation \(ed\equiv1\pmod{\varphi(n)}\).}
\end{codepairs}

Now encryption raises the message to \(e\) modulo \(n\), and decryption raises the
ciphertext to \(d\):

\begin{codepairs}
\pair{number-theory/rsa/11}{The plaintext number \(42\).}
\pair{number-theory/rsa/12}{Encrypted, it becomes \(2557\)---bearing no visible relation
to the original.}
\pair{number-theory/rsa/13}{Decryption with the private exponent returns \(42\) exactly.}
\pair{number-theory/rsa/14}{The round trip closes.}
\end{codepairs}

Why it works is Euler's theorem from \S\ref{ssec:nt-modular}: because \(ed\equiv1\pmod{\varphi(n)}\),
raising the message to \(e\) and then to \(d\) raises it to a power \(\equiv1\), which
leaves it unchanged modulo \(n\). Anyone holding \((n,e)\) can encrypt; only the holder of
\(d\)---recoverable only by factoring \(n\)---can decrypt. Every step ran on integers no
larger than a few thousand, but nothing about the method changes when the primes have a
hundred digits each; only the factoring stays out of reach.

\subsection*{Where this connects}

Number theory is the exact spine of Mathilda's arithmetic. Its foundation is the integer
model of \S\ref{sec:arithmetic}---the same floored \B{Mod} and GMP big integers, here put
to work. Its natural sequel is the arithmetic of \emph{solutions}: the classical Diophantine
questions---Pell's equation, Pythagorean triples, sums of squares, and the rest---are
answered by \B{Solve} and \B{Reduce} over the integers, treated in
\S\ref{ssec:alg-diophantine}, where the continued fractions of \S\ref{ssec:nt-contfrac}
reappear as the engine behind Pell. The polynomial \B{Factor} of \S\ref{sec:algebra} is the
algebraic mirror of \B{FactorInteger}, square-free decomposition and Hensel lifting
standing in for trial division and rho. And the integer-valued special functions that shade
number theory into analysis---\B{BernoulliB}, \B{Zeta}, and their kin---await in
\S\ref{sec:specialfns}, including the Dirichlet characters and \(L\)-functions of
\S\ref{ssec:sf-dirichlet}, where the residue classes of \S\ref{ssec:nt-modular} become
analytic objects. Underneath every result here is GMP, and the same exact integers
feed the compiler of Chapter~\ref{ch:compilation} when a number-theoretic loop needs to run
at machine speed.
